\begin{center}
  \centering
  \footnotesize
  \setlength{\tabcolsep}{4pt}
  \medskip\noindent\begin{minipage}{\textwidth}\centering
  \textbf{Status}\par\smallskip
  \begin{tabularx}{\textwidth}{@{}
      >{\raggedright\arraybackslash}p{0.24\textwidth}
      >{\raggedright\arraybackslash}X
      >{\raggedright\arraybackslash}p{0.10\textwidth}@{}}
    \toprule
    Name & One-line claim & Substrate cite \\
    \midrule
    Objectivity as Common Quotient (\Fref{F17}) &
    For surjective access and public quotients, when the public quotient is the common quotient of the accesses (it factors through each, and every map factoring through each factors through it), a readout factors through every access exactly when the public quotient determines it, and exactly when no access identifies two histories with different readout values.  &
    \textemdash \\
    Lawhood-Descent Theorem (\Fref{F18}) &
    For a surjective access quotient, a candidate law is lawful exactly
    when it factors through it.  &
    \textemdash \\
    Interface-Mediation Theorem (\Fref{F22}) &
    For surjective source quotient and boundary readout, an update is mediated iff its target class descends, uniquely, through (source
    class, boundary value); failure splits into boundary leakage and internal
    insufficiency.  &
    \textemdash \\
    Probability as Stable Unresolved Fiber (\Fref{F23}) &
    For a surjective quotient, a readout either descends or is unresolved, never both; stable
    probability is unresolvedness together with the declared ledger
    conditions Stable, which do not involve the weights and hold trivially for empty supports and a one-point target set; on a finite carrier, with the weight assignment given on fiber events by real fiber
    weights, positive and summing to one on each fiber, descent holds iff every fiber chance is a point mass (value 1 at one readout value).  &
    \textemdash \\
    \bottomrule
  \end{tabularx}
  \end{minipage}\par

  \medskip\noindent\begin{minipage}{\textwidth}\centering
  \textbf{Records}\par\smallskip
  \begin{tabularx}{\textwidth}{@{}
      >{\raggedright\arraybackslash}p{0.24\textwidth}
      >{\raggedright\arraybackslash}X
      >{\raggedright\arraybackslash}p{0.10\textwidth}@{}}
    \toprule
    Name & One-line claim & Substrate cite \\
    \midrule
    Object-Persistence Theorem (\Fref{F19}) &
    For a surjective source access quotient \(q_i\), an object persists exactly when its split obstruction is empty; exactly one of four comparison statuses holds; the pairing \(h\mapsto(q_i(h),q_j(\gamma h))\) retains \(q_i\) and carries \(q_j\circ\gamma\); viability, relations and holonomy hold exactly when their violation sets are empty.  &
    \textemdash \\
    Fact--Record Theorem (\Fref{F20}) &
    For surjective record signatures \(s\) and \(\rho_s\), an event is recorded, a record is stable, and a fact value is stable
    exactly when the corresponding obstruction is empty; the record
    comparison has exactly one of four statuses.  &
    \textemdash \\
    Reflexive-Nonclosure Theorem (\Fref{F21}) &
    Under rules 2--4 of a same-layer audit policy, each of which is needed, a complete self-soundness claim is never accepted, and, when the meta layer accepts no self-sound meta-claim, a meta-layer certificate comes from an accepted
    meta-claim that is not itself self-sound.  &
    \textemdash \\
    Causality Theorem (\Fref{F25}) &
    For surjective context and intervention quotients, intervention
    transport descends to (context class, intervention
    class) exactly when its obstruction is empty, and then uniquely;
    under a descended transport every stable causal effect has a witness
    at one raw context with two raw interventions.  &
    \textemdash \\
    \bottomrule
  \end{tabularx}
  \end{minipage}\par

  \medskip\noindent\begin{minipage}{\textwidth}\centering
  \textbf{Coherence}\par\smallskip
  \begin{tabularx}{\textwidth}{@{}
      >{\raggedright\arraybackslash}p{0.24\textwidth}
      >{\raggedright\arraybackslash}X
      >{\raggedright\arraybackslash}p{0.10\textwidth}@{}}
    \toprule
    Name & One-line claim & Substrate cite \\
    \midrule
    Moduli--Contingency Theorem (\Fref{F26}) &
    When the moduli map identifies exactly the equivalent lawful closures and every modulus is attained by a lawful closure, a readout agrees
    on lawful closures with a map on the moduli set
    exactly when it is presentation-independent on lawful closures; it is then either
    necessary or contingent, never both.  &
    \textemdash \\
    Conservation Theorem (\Fref{F27}) &
    Conserved quantities are exactly readouts that factor through
    route-orbit quotients.  &
    \textemdash \\
    Composability Theorem (\Fref{F28}) &
    For a surjective part/interface quotient \(\kappa\), the composite outcome, cross residual and admissibility each descend through \(\kappa\), uniquely, exactly when their split-pair sets are empty; exact composability is emptiness of all three.  &
    \textemdash \\
    Presentation-Invariance Theorem (\Fref{F29}) &
    For a quotient with a section, a readout constant on presentation classes is exactly one that factors through the quotient, agrees with its value at the section representative, or is constant along any relation generating the kernel; relation invariance under any tuple quotient is descent; presentation invariance does not imply equivariance.  &
    \textemdash \\
    Explanation-as-Compression Theorem (\Fref{F30}) &
    For surjective base and explanation readouts, an explanation is exact
    exactly when its source and target obstructions are empty, and is a
    strict compression exactly when it is determined by the base and merges
    two histories the base separates; a dynamic explanation makes the later
    target descend through the explanation.  &
    \textemdash \\
    Counterfactual-Legitimacy Theorem (\Fref{F31}) &
    A counterfactual is legitimate exactly when both obstruction sets are empty, exactly when admissibility, and target truth on admissible pairs, descend to (context class, intervention class).  &
    \textemdash \\
    \bottomrule
  \end{tabularx}
  \end{minipage}\par

  \captionof{table}{Status, Records, Coherence at a glance: the fourteen laws of \Cref{sec:status-records-coherence}, grouped into reading groups.}
  \label{tab:axis-status-records-coherence-summary}
\end{center}
